% !TeX root = ../drafts/07-instruction-tables.tex
\section{The instruction tables}\label{sec:tables}

Each instruction has its own table. Throughout, an operand address is written $\loc{o}=\fp\cdot o$ and is a bus coordinate, not a column: no table commits one, and none constrains one, since a coordinate may be a product of two columns (\S\ref{sec:m3}).

\subsection{\texttt{XOR}}\label{sec:tab-xor}

Asserts $\loc{o_C}=\loc{o_A}+\loc{o_B}$ in $\K$ (bitwise \textsc{xor}).
\begin{itemize}
\item \textbf{Columns:} $\pc,\fp$; operands $o_A,o_B,o_C$; addresses $a_A,a_B,a_C$; values $v_A,v_B,v_C\in\K$; memory counts $r_A,r_B,r_C$; bytecode count $r_{\mathrm{bc}}$.
\item \textbf{Constraints:} the three addresses $a_X=\fp\cdot o_X$ and the base-field sum $v_C=v_A+v_B$.
\item \textbf{Flushes:}
  \begin{itemize}
  \item \textbf{State:} \pull\ $\tup{\dsep{ST},\pc,\fp}$, \push\ $\tup{\dsep{ST},\gen\cdot\pc,\fp}$.
  \item \textbf{Bytecode:} \pull\ $\tup{\dsep{BC},\pc,r_{\mathrm{bc}},\opc{XOR},o_A,o_B,o_C,0,0,0,0,0}$, \push\ the same at $\gen\cdot r_{\mathrm{bc}}$.
  \item \textbf{Memory} ($X\in\{A,B,C\}$): \pull\ $\tup{\dsep{MEM},a_X,r_X,v_X}$, \push\ $\tup{\dsep{MEM},a_X,\gen\cdot r_X,v_X}$.
  \end{itemize}
\end{itemize}

\subsection{\texttt{MUL}}\label{sec:tab-mul}

Identical to \texttt{XOR} but for the base-field product:
\begin{itemize}
\item \textbf{Columns:} as \texttt{XOR}.
\item \textbf{Constraints:} the three addresses and $v_C=v_Av_B$ in $\K$.
\item \textbf{Flushes:} as \texttt{XOR}, with opcode $\opc{MUL}$.
\end{itemize}

\subsection{\texttt{XOR\_192} and \texttt{MUL\_192}}\label{sec:tab-ext}

Each operand is three consecutive base words. The table commits one address
$a_X=\fp o_X$ per operand and uses $\gen a_X,\gen^2a_X$ virtually for the
remaining coordinates. It performs nine memory accesses, reassembles
$v_X=v_{X,0}+v_{X,1}y+v_{X,2}y^2$, and constrains respectively
$v_C=v_A+v_B$ or $v_C=v_Av_B$ in $\E$. The multiplication table also has a
public-bytecode mode $m=1$ for $v_A\in\K$: only $v_{A,0}$ participates in the
product, while the other two memory cells remain ordinary read-only memory-bus
accesses. Two quadratic constraints tie the effective upper limbs to the
memory values when $m=0$ and to zero when $m=1$, so the AIR remains degree two.
In the zkDSL, \texttt{xor\_192} exposes addition,
\texttt{mul\_192\_base} exposes the scalar mode, and \texttt{div\_192} is a
checked multiplication with the quotient back-solved by witness generation.

\subsection{\texttt{SET\_CONSTANT}}\label{sec:tab-set}

Asserts $\loc{o}=k\in\K$, how a constant enters the read-only memory.
\begin{itemize}
\item \textbf{Columns:} $\pc,\fp$; operand $o$; immediate $k\in\K$; address $a$; memory count $r$; bytecode count $r_{\mathrm{bc}}$.
\item \textbf{Constraints:} the address $a=\fp\cdot o$.
\item \textbf{Flushes:}
  \begin{itemize}
  \item \textbf{State:} \pull\ $\tup{\dsep{ST},\pc,\fp}$, \push\ $\tup{\dsep{ST},\gen\cdot\pc,\fp}$.
  \item \textbf{Bytecode:} \pull\ $\tup{\dsep{BC},\pc,r_{\mathrm{bc}},\opc{SET},o,k,0,0,0,0,0,0}$, \push\ the same at $\gen\cdot r_{\mathrm{bc}}$.
  \item \textbf{Memory:} \pull\ $\tup{\dsep{MEM},a,r,k}$, \push\ $\tup{\dsep{MEM},a,\gen\cdot r,k}$.
  \end{itemize}
\end{itemize}

\subsection{\texttt{DEREF}}\label{sec:tab-deref}

\texttt{DEREF} follows a pointer and asserts the cell it points to. The row reads three cells:
\begin{itemize}
\item the \emph{pointer} $p$, at $\fp\cdot\alpha$;
\item the \emph{target} $v_2$, at the dereferenced address $p\cdot\beta$;
\item the \emph{local} value $v_3$, at $\fp\cdot\gamma$.
\end{itemize}
The store asserts that $v_2$ equals the source chosen by the mode $s$ (\S\ref{sec:vm}): the local cell $v_3$, the return address $\gen^{2}\cdot\pc$ (a virtual $\times\gen^{2}$ of $\pc$, \S\ref{sec:prelim}), or the frame pointer $\fp$. The mode is two boolean flags $(f_{\pc},f_{\fp})$: $(0,0)$ for $\texttt{deref\_cell}[\gamma]$, $(1,0)$ for $\texttt{deref\_pc}$, $(0,1)$ for $\texttt{deref\_fp}$.
\begin{itemize}
\item \textbf{Columns:} $\pc,\fp$; operands $\alpha,\beta,\gamma$; flags $f_{\pc},f_{\fp}$; addresses $a_1,a_2,a_3$; words $p,v_2,v_3\in\K$; memory counts $r_1,r_2,r_3$; bytecode count $r_{\mathrm{bc}}$.
\item \textbf{Constraints:} the addresses $a_1=\fp\cdot\alpha$, $a_2=p\cdot\beta$, $a_3=\fp\cdot\gamma$, and the store
\[
  v_2 \;=\; (1+f_{\pc}+f_{\fp})\,v_3 \;+\; f_{\pc}\,(\gen^{2}\cdot\pc) \;+\; f_{\fp}\,\fp ,
\]
which gives $v_3$, $\gen^{2}\cdot\pc$, or $\fp$ at the three flag settings. All relations are in $\K$. The flags need no validity check: the verifier reads them from the public program (\S\ref{sec:e2e-bc}).
\item \textbf{Flushes:}
  \begin{itemize}
  \item \textbf{State:} \pull\ $\tup{\dsep{ST},\pc,\fp}$, \push\ $\tup{\dsep{ST},\gen\cdot\pc,\fp}$.
  \item \textbf{Bytecode:} \pull\ $\tup{\dsep{BC},\pc,r_{\mathrm{bc}},\opc{DRF},\alpha,\beta,\gamma,f_{\pc},f_{\fp},0,0,0}$, \push\ the same at $\gen\cdot r_{\mathrm{bc}}$.
  \item \textbf{Memory:} three reads, each \pull\ at count $r$ and \push\ at $\gen\cdot r$: pointer $\tup{\dsep{MEM},a_1,r_1,p}$, target $\tup{\dsep{MEM},a_2,r_2,v_2}$, local cell $\tup{\dsep{MEM},a_3,r_3,v_3}$.
  \end{itemize}
\end{itemize}

\subsection{\texttt{DEREF\_WIDE}}\label{sec:tab-deref-ext}

\texttt{DEREF\_WIDE} bundles consecutive scalar cell-mode dereferences. It
reads the pointer $p$ at $a_1=\fp\alpha$, the heap run beginning at
$a_2=p\beta$, and the local run beginning at $a_3=\fp\gamma$. The successors
$\gen a_2,\gen^2a_2$ and $\gen a_3,\gen^2a_3$ are virtual columns. A width bit
$w$, fixed by public bytecode, selects native two-word \texttt{DEREF\_128}
($w=0$) or three-word \texttt{DEREF\_192} ($w=1$). Its constraints are the
three base-address relations and the packed equality
\[
  v_{2,0}+yv_{2,1}+wy^2v_{2,2}
  =v_{3,0}+yv_{3,1}+wy^2v_{3,2}.
\]
The row always flushes one pointer access and all six run-cell accesses to
memory, plus the usual state and bytecode interactions. In two-word mode the
third cells are independently memory-bound padding reads, but are not equated.
Thus one VM cycle transfers either native 128-bit data or one three-word
extension value without weakening the 64-bit memory-word invariant.

\subsection{\texttt{JUMP}}\label{sec:tab-jump}

A conditional jump on whether $c=\loc{o_c}\in\K$ is nonzero. It transfers to $\pc\gets d$, $\fp\gets f$ when $c\neq0$ and falls through otherwise. A committed boolean indicator $b=[\,c\neq0\,]$ and inverse $w\in\K$ certify
\[
  b=c\cdot w,\qquad c\,(b+1)=0 .
\]
When $c\neq0$ the second forces $b=1$ and the first then $w=c^{-1}$; when $c=0$ the first forces $b=0$. Writing $s=\gen\cdot\pc$ for the virtual fall-through, the committed next state is the selection
\[
  \mathit{next\_pc}=b\cdot d+(b+1)\,s,\qquad \mathit{next\_fp}=b\cdot f+(b+1)\,\fp.
\]
Alone among the tables, \texttt{JUMP} commits its successors.
\begin{itemize}
\item \textbf{Columns:} $\pc,\fp,\mathit{next\_pc},\mathit{next\_fp}$; operands $o_c,o_d,o_f$; addresses $a_c,a_d,a_f$; words $c,d,f,w\in\K$; indicator $b\in\{0,1\}$; memory counts $r_c,r_d,r_f$; bytecode count $r_{\mathrm{bc}}$.
\item \textbf{Constraints:} the three addresses, the indicator pair $b=cw$ and $c(b+1)=0$, and the two next-state selections above, all over $\K$.
\item \textbf{Flushes:}
  \begin{itemize}
  \item \textbf{State:} \pull\ $\tup{\dsep{ST},\pc,\fp}$, \push\ $\tup{\dsep{ST},\mathit{next\_pc},\mathit{next\_fp}}$.
  \item \textbf{Bytecode:} \pull\ $\tup{\dsep{BC},\pc,r_{\mathrm{bc}},\opc{JMP},o_c,o_d,o_f,0,0,0,0,0}$, \push\ the same at $\gen\cdot r_{\mathrm{bc}}$.
  \item \textbf{Memory} ($X\in\{c,d,f\}$): \pull\ $\tup{\dsep{MEM},a_X,r_X,X}$, \push\ $\tup{\dsep{MEM},a_X,\gen\cdot r_X,X}$.
  \end{itemize}
\end{itemize}


\subsection{\texttt{BLAKE3}}\label{sec:tab-blake3}

Compresses two 256-bit operands to a 256-bit digest. Each input is two independently addressed 128-bit chunks, based at $\fp\cdot o_{A,i}$ and $\fp\cdot o_{B,i}$ for $i\in\{0,1\}$; the second word of each chunk is the same product times $\gen$. The four-word chaining value and output are based at $\fp\cdot o_{CV}$ and $\fp\cdot o_C$, their successors the same products times $\gen,\gen^2,\gen^3$. Each row therefore accesses sixteen memory cells while committing no address at all: the $\gen^k$ factor on a bus product reaches a run's $k$-th successor for free.

Flock~\cite{flock-impl} proves the compression relation with a batched R1CS over a Boolean witness. The witness is packed into $q_\flock$, with $64$ bits per $\K$-element. $q_\flock$ is the one aligned exception to the jagged layout (\S\ref{sec:stacking}): it occupies the offset-zero prefix, so its ring-switch and fixed-slot claims keep their direct subcube formulas. The concrete ring switch of \S\ref{sec:ringswitch} maps its $\F_2$ claims to one $\E$-weighted claim on $q_\flock$, which joins the batched opening (\S\ref{sec:e2e-unrolled}).

The eighteen $64$-bit words of $a,b,c,cv,m_0,m_1$ are contained in $q_\flock$ and need no separate value-column commitment. Each is a strided point evaluation of $q_\flock$ and is linked directly to its single-word memory interaction. The R1CS matrices prove the compression relation for those supplied values.
\begin{itemize}
\item \textbf{Columns:} $\pc,\fp$; operands $o_{A,0},o_{A,1},o_{B,0},o_{B,1},o_{CV},o_C$ and metadata $m_0,m_1$; their six starting addresses; sixteen per-word memory counts; bytecode count $r_{\mathrm{bc}}$. The eighteen values live in $q_\flock$.
\item \textbf{Constraints:} the six address relations $a_X=\fp o_X$. Flock proves $c=\mathrm{BLAKE3}(a,b;cv,m_0,m_1)$.
\item \textbf{Flushes:}
  \begin{itemize}
  \item \textbf{State:} \pull\ $\tup{\dsep{ST},\pc,\fp}$, \push\ $\tup{\dsep{ST},\gen\cdot\pc,\fp}$.
  \item \textbf{Bytecode:} \pull\ $\tup{\dsep{BC},\pc,r_{\mathrm{bc}},\opc{B3},o_{A,0},o_{A,1},o_{B,0},o_{B,1},o_{CV},o_C,m_0,m_1}$, \push\ the same at $\gen\cdot r_{\mathrm{bc}}$.
  \item \textbf{Memory:} sixteen reads: two at $a_{X,i},\gen a_{X,i}$ for each input chunk, and four at $\gen^j a_{CV}$ and $\gen^j a_C$ for $j\in\{0,1,2,3\}$, each paired with the count multiplied by $\gen$.
  \end{itemize}
\end{itemize}
